%% Distinct first-order calculi proofs retained after exact source restoration.
\noindent\textit{Reader note.} These two separately attributed editorial proofs
support only the retained equality and Hilbert additions.  The inherited
reduct and soundness proofs remain exact in the body and are not repeated;
historical items .1--.3 remain on disk only.  Registry sheets provide
navigation and variant context, not source replacement.

\begin{annexitem}{annex:p2.fol.calc.4}{thm:p02-equality-laws}{Derived equality laws}
\begin{proof}
Reflexivity and formula replacement are the two registered equality rules.
Choose \(z\notin
\operatorname{Var}(r)\cup\operatorname{Var}(t)\cup\operatorname{Var}(u)\).
For symmetry, take \(r=r\) and apply replacement to \(r=t\) with
\(\varphi(z)\defeq z=r\), obtaining \(t=r\).  For transitivity, apply
replacement to \(t=u\) and \(r=t\) with
\(\varphi(z)\defeq r=z\), obtaining \(r=u\).  Total capture-avoiding
substitution makes both displayed instances legitimate.
\end{proof}
\end{annexitem}

\begin{annexitem}{annex:p2.fol.calc.5}{thm:p02-hilbert-fol-deduction}{Hilbert-FOL deduction theorem}
\begin{proof}
The reverse implication appends \(\varphi\) and applies modus ponens.  For
the forward implication, induct through a dependency-annotated derivation
from \(\Gamma\cup\{\varphi\}\), proving
\(\Gamma\derives_{\mathsf H_{\mathsf{FOL}}}\varphi\to\chi_i\).
Premises from \(\Gamma\) and axioms use A1; the line \(\varphi\) uses the
standard A1/A2 derivation of \(\varphi\to\varphi\); modus ponens uses A2.

Suppose \(\chi_i=\forall x\,\psi\) comes by generalization.  If the
dependency set for \(\psi\) contains \(\varphi\), the rule condition gives
\(x\notin\FV(\varphi)\).  Generalize the induction result
\(\varphi\to\psi\), then use the \(\forall2\) schema to obtain
  \(\varphi\to\forall x\psi\).  If \(\varphi\) is absent from that dependency
  set, its traced subderivation is already from \(\Gamma\); generalize
  \(\psi\) with the same dependencies and use A1 to infer
  \(\varphi\to\forall x\psi\).

  Suppose instead that \(\chi_i=\exists x\,\theta\to\psi\) follows by
  existential elimination from \((\theta\to\psi,D)\).  If
  \(\varphi\notin D\), the traced subderivation is already from \(\Gamma\);
  apply existential elimination there and then A1.  If \(\varphi\in D\), the
  rule condition gives \(x\notin\FV(\varphi)\) as well as
  \(x\notin\FV(\psi)\).  From the induction result
  \(\varphi\to(\theta\to\psi)\), an A1--A3 propositional permutation gives
  \(\theta\to(\varphi\to\psi)\).  Existential elimination yields
  \(\exists x\,\theta\to(\varphi\to\psi)\), and the inverse permutation gives
  \(\varphi\to(\exists x\,\theta\to\psi)\).

  Thus every annotated line is covered, and the
  last one proves the claim without unrestricted generalization.
\end{proof}
\end{annexitem}
